\documentclass[conference]{IEEEtran}
\IEEEoverridecommandlockouts
% The preceding line is only needed to identify funding in the first footnote. If that is unneeded, please comment it out.
\usepackage{cite}
\usepackage{amsmath,amssymb,amsfonts}
\usepackage{algorithmic}
\usepackage{graphicx}
\usepackage{textcomp}
\usepackage{xcolor}
\def\BibTeX{{\rm B\kern-.05em{\sc i\kern-.025em b}\kern-.08em
    T\kern-.1667em\lower.7ex\hbox{E}\kern-.125emX}}
\begin{document}

\title{Dư đoán giá cổ phiếu dựa trên Bayesian-LSTM\\}
\author{\IEEEauthorblockN{1\textsuperscript{st} Phạm Hùng Vỹ}
\IEEEauthorblockA{\textit{dept. name of organization (of Aff.)} \\
\textit{name of organization (of Aff.)}\\
City, Country \\
email address or ORCID}
\and
\IEEEauthorblockN{2\textsuperscript{nd} Phan Hữu Chí}
\IEEEauthorblockA{\textit{dept. name of organization (of Aff.)} \\
\textit{name of organization (of Aff.)}\\7
City, Country \\
email address or ORCID}
}

\maketitle

\begin{abstract}
In this research, we study the problem of stock market forecasting using Recurrent
Neural Network(RNN) with Long Short-Term Memory (LSTM). The purpose of this
research is to examine the feasibility and performance of LSTM in stock market
forecasting. We optimize the LSTM model by testing different configurations, i.e., the
number of neurons in hidden layers and number of samples in sequence. Instead of
using daily stock price data, we collect hourly stock data from the IQFEED database in
order to train our model with relatively low noise samples. Nevertheless, based on the
prediction results of LSTM model, we build up a stock database with six U.S market
stocks from five different industries. The average test accuracy of these six stocks is
54.83%, where the highest accuracy is at 59.5% while the lowest is at 49.75%. We then
develop a trade simulator to evaluate the performance of our model by investing the
portfolio within a period of 400 hours, the total profit gained by the model is
$413,233.33 with $6,000,000 initial investment.
\end{abstract}
\begin{IEEEkeywords}
Thị trường chứng khoán; B-LSTM; Dư đoán giá chứng khoán;Tối ưu hoá Bayesian
\end{IEEEkeywords}

\section{Giới thiệu}
Thị trường chứng khoán là các tổ chức giao dịch nơi chứng khoán (vốn chủ sở hữu) và tài chính khác
các công cụ như trái phiếu được cung cấp cho thương mại. Đối với cổ phiếu, thị trường thường hoạt động
một giao dịch người mua đưa ra mong muốn mức giá muốn mua, người bán đưa ra mức giá muốn bán, và nếu có người mua và người bán đều có mức giá phù hợp với nhau thì giao dịch sẽ được diễn ra. Nếu không thì sẽ không có giao dịch nào diễn ra và chờ đợi một mức giá  trong tương lai hoặc hết hạn.

Trong hầu hết các sàn giao dịch chứng khoán, thị trường phổ biến và dễ tiếp cận là thị trường chứng khoán
(cổ phiếu), có rất ít rào cảng để mọi người có thể tham gia. Thị trường chứng khoán là
do đó tích cực hơn, có nhiều người chơi và do đó một phân khúc xứng đáng để nghiên cứu thêm.
Hiệu suất của thị trường chứng khoán được đo lường hàng ngày bởi một số chỉ số chính
chẳng hạn như 'chỉ số tổng hợp', chỉ số thị trường chứng khoán của tất cả các cổ phiếu được giao dịch tại Sở giao dịch chứng khoán. Một chỉ số như vậy rất quan trọng trong việc không chỉ đo lường
hiệu suất của các giao dịch trên thị trường chứng khoán mà còn phản ánh rõ các hoạt động kinh tế của một quốc gia và các mối quan hệ quốc tế. Cổ đông tuy nhiên không trực tiếp thực hiện giao dịch,
cũng không có cuộc họp nào giữa người mua và người bán để đàm phán. Cổ đông giao dịch
bằng cách đưa ra hướng dẫn cho các Môi giới chứng khoán của họ, những người lần lượt thực hiện các lệnh. Môi giới chứng khoán
thường cũng tư vấn cho khách hàng về nơi giao dịch. Trong vai trò tư vấn của họ, một số môi giới chứng khoán
căn cứ lời khuyên của họ về các nguyên tắc cơ bản của các cổ phiếu khác nhau hoặc thực hiện kỹ thuật
phân tích. Tuy nhiên, không có phương pháp dự đoán nào trong số này đảm bảo lợi nhuận vì chúng
thường chỉ cho thấy một xu hướng trong tương lai và khả năng tăng hoặc giảm giá chứ không phải là
giá cổ phiếu dự kiến ​​trong tương lai. Môi giới chứng khoán cần được trao quyền, thông qua tốt hơn
công cụ dự đoán, để cho phép họ có một số khả năng để cung cấp lời khuyên tốt nhất cho họ
khách hàng.Mọi người người môi giới chứng khoán đều mong muốn có một công cụ dự đoán mà có thể sử dụng để phán đoán về biến động giá chính xác như là một cơ sở của đầu tư. 
Đây là công cụ mọi người đều mong muốn có kể từ khi thị trường chứng khoán ra đời. Các nhà toán học từ những năm 1960 đã đưa ra các mô hình thống kê như ARIMA, SARIMA để nắm bắt thị trường nhưng kết quả đạt được lại không như mong đợi.
Kể từ năm 2010, máy học đã đạt đươc những đôt phá nhất định. Ý tưởng chinh phục cổ phiếu đã quay trở lại và mãnh liệt hơn khi học sâu đã vượt qua hàng loạt các thuật toán khác chứng tỏ sự yêu việc của mình. Đặc biệt khi sử dụng các mạng thần kinh tuần tự cho các vấn đề chuỗi thời gian cho kết quả chính xác ngoài mong đợi
Nhưng việc chọn mô hình và đưa ra các siêu tham số là một công viêc cần có kinh nghiệm chuyên sâu về cả cổ phiếu lẫn máy học. 
Không phải ai cũng có khả năng đó. Vì vậy chúng tối đễ xuất tối sử dụng tối ưu hoá Bayesian và tìm cách xây dựng mô hình sử dụng LSTM dễ tạo các mô hình có nhiều siêu tham số để tối ưu để mô hình đưa ra


\section{Các nghiên cứu liên quan}
\subsection{Sử dụng máy học trong dự đoán chứng khoán}
Dự báo giá cổ phiếu là một việc rất phức tạp. Hầu hết các nhà môi giới chứng khoán sử dụng phân tích chuỗi kỹ thuật, cơ bản hoặc phân tích chuỗi thời gian (time series) trong việc cố gắng dự đoán giá cổ phiếu. Tuy nhiên, các chiến lược này không dẫn đến kết quả đáng tin cậy vì chúng hướng dẫn về xu hướng và không phải là giá có độ chính xác cao nhất. Cần phải sử dụng các phương pháp nâng cao để dự đoán kết quả chính xác nhất. Các nhà nghiên cứu đã sử dụng các phương pháp khác nhau và các bộ tham số đầu vào khác nhau để dự đoán giá cổ phiếu trong vài thập kỷ qua.

Adebiyi và Adewumi [3] đã xây dựng mô hình dự đoán giá cổ phiếu bằng cách sử dụng mô hình Autoregressive integrated moving average (ARIMA). Kết quả cho thấy mô hình ARIMA có tiềm năng mạnh mẽ để dự đoán ngắn hạn.

Xing, Sun, Wang và Yu. [24] giới thiệu một loại phương pháp dựa trên Mô hình Markov ẩn để dự báo xu hướng giá cổ phiếu. Khác với mô hình dự đoán cổ phiếu hiện tại, họ đã cố gắng tìm mối quan hệ ẩn tồn tại giữa các giá cổ phiếu bằng Mô hình Markov ẩn. Kết quả thử nghiệm cho thấy, phương pháp này có thể nhận được kết quả khá chính xác, đặc biệt hiệu quả trong dự đoán thời gian ngắn.

Marček [4] đã mô tả khái niệm cơ bản của mô hình hồi quy tuyến tính mờ dựa trên nguyên tắc mở rộng tham số mờ. Bài viết đã trình bày mô hình tự phát (AR) sử dụng nguyên tắc mở rộng tham số mờ và Fuzzy Neural Network (FNN) để ước tính và dự đoán giá cổ phiếu. Qua đó cho thấy rằng kết quả ban đầu từ việc sử dụng kiến trúc FNN tốt hơn kiến trúc Artificial Neural Network (ANN) căn bản cho việc dự đoán hằng ngày.

Hegazy, Soliman và Salam [5] đã triển khai mô hình học máy để dự đoán giá thị trường chứng khoán. Bằng cách kết hợp thuật toán Particle swarm optimization(PSO) và thuật toán Least square support vector machine (LS-SVM). Thuật toán PSO chọn ba tham số bất kì kết hợp tốt nhất từ nghiên cứu về dữ liệu lịch sử của cổ phiếu và các chỉ số kỹ thuật cho LS-SVM để tránh các vấn đề over-fitting. Mô hình được đề xuất đã được áp dụng và đánh giá bằng cách sử dụng mười ba bộ dữ liệu tài chính và so sánh với mạng nơ ron nhân tạo với thuật toán Levenberg - Marquest (LM). Kết quả thu được cho thấy mô hình đề xuất có độ chính xác dự đoán tốt hơn và tiềm năng của thuật toán PSO trong việc tối ưu hóa LS-SVM.

Tiwari, Bharadwaj và Gupta [22] đề xuất mục đích của việc phân tích dữ liệu được sử dụng để hỗ trợ các nhà đầu tư đưa ra dự đoán tài chính chính xác để đưa ra các quyết định đầu tư đúng đắn. Hai nền tảng đã được sử dụng: Python và R

Các kỹ thuật khác nhau như Arima, Holt winters, mạng nơ-ron (Feed forward và Multi-layer perceptron), hồi quy tuyến tính và time series được triển khai để dự báo hiệu suất giá chỉ số mở trong R. Mặc dù trong Python Multi-layer perceptron và hồi quy vectơ đã được triển khai để dự báo giá cổ phiếu Nifty 50 và phân tích cá nhân của cổ phiếu đã được thực hiện bằng cách sử dụng các tweet gần nhất trên Twitter. Các chỉ số chứng khoán Nifty 50 (ANSEI) được coi là đầu vào dữ liệu cho các phương thức được triển khai. Chín năm dữ liệu đã được sử dụng. Độ chính xác được tính bằng cách sử dụng 2-3 năm kết quả dự báo của R và 2 tháng kết quả dự báo của Python sau khi so sánh với giá thực tế của các cổ phiếu.

Mittal và Goel [6] đã áp dụng các nguyên tắc phân tích cá nhân và machine learning để tìm ra mối tương quan giữa ‘public sentiment’ and ‘market sentiment’. Dữ liệu Twitter được sử dụng để dự đoán tâm trạng công cộng và những ngày trước đó Giá trị trung bình công nghiệp (DJIA) của Dow Dow Jones đã được sử dụng để dự đoán diễn biến của thị trường chứng khoán. Họ đề xuất một phương pháp xác thực chéo mới cho dữ liệu tài chính để kiểm tra kết quả và thu được độ chính xác 75,56\% bằng cách sử dụng Self Organizing Fuzzy Neural Networks (SOFNN) trên các nguồn cấp dữ liệu Twitter và DJIA từ giai đoạn tháng 6 năm 2009 đến tháng 12 năm 2009.

Wamkaya và Lawrence [7] đã đề xuất việc sử dụng Mạng nơ ron nhân tạo là feedforward multi-layer perceptron with error backpropagation (Ở trên thì nó viết là feedforward and multi-layer perceptron). Họ đã phát triển một mô hình cấu hình 5: 21: 21: 1 với 80\% dữ liệu đào tạo trong 130.000 chu kỳ. Nghiên cứu đã phát triển một nguyên mẫu và thử nghiệm nó trên dữ liệu 2008 - 2012 từ các thị trường chứng khoán nơi kết quả dự đoán cho thấy MAPE trong khoảng 0,71\% đến 2,77\%.

Dunne [27] đã phân tích các phương pháp dự đoán thị trường chứng khoán hiện có và các phương pháp mới. Ông đã thực hiện ba cách tiếp cận khác nhau : Phân tích cơ bản, Phân tích kỹ thuật và ứng dụng Machine Learning. Ông tìm thấy bằng chứng ủng hộ mô hình yếu của Efficient Market Hypothesis, rằng giá lịch sử không chứa thông tin hữu ích nhưng ngoài dữ liệu mẫu có thể là dự đoán. Ông đã chỉ ra rằng Fundamental Analysis và Machine Learning có thể được sử dụng để hướng dẫn các nhà đầu tư đưa ra quyết định. Ông chứng minh một lỗ hổng phổ biến trong phương pháp Phân tích Kỹ thuật và cho thấy rằng nó tạo ra thông tin hữu ích nhưng hạn chế. Dựa trên những phát hiện của ông, các chương trình giao dịch thuật toán đã được phát triển và mô phỏng bằng Quantopian.

Selvin, Vinayakumar, Gopalakrishnan, Menonans và Soman K.P [8] đã thử nghiệm một mô hình tiếp cận độc lập. Thay vì lắp dữ liệu vào một mô hình cụ thể, họ đã xác định động lực tiềm ẩn hiện có trong dữ liệu bằng kiến trúc deep learning. Họ đã sử dụng ba kiến trúc deep learning khác nhau để dự đoán giá của các công ty niêm yết NSE và so sánh hiệu suất của họ. Ba mô hình là NN, LSTM và CNN. Hiệu suất của các mô hình đã được định lượng bằng lỗi phần trăm. Giá trị tối đa của tỷ lệ phần trăm lỗi đã thu được cho mỗi mô hình. Kết quả cho thấy CNN đã cho kết quả chính xác hơn so với hai mô hình còn lại.

Kita, Zuo, Harada và Mizuno [9] đã phát triển thuật toán dự đoán giá cổ phiếu bằng cách sử dụng mạng Bayes. Thuật toán sử dụng mạng hai lần. Đầu tiên, mạng được xác định từ giá cổ phiếu hàng ngày và sau đó nó được áp dụng để dự đoán giá cổ phiếu hàng ngày đã được quan sát. Các lỗi dự đoán được đánh giá từ giá cổ phiếu hàng ngày và dự đoán của nó. Thứ hai, mạng được xác định một lần nữa từ cả giá cổ phiếu hàng ngày và lỗi dự đoán hàng ngày và sau đó nó được áp dụng cho dự đoán giá cổ phiếu trong tương lai. Kết quả bằng số cho thấy lỗi dự đoán tối đa của thuật toán hiện tại là 30\%.


Xiongwen Pang, Yanqiang Zhou, Pan Wang, Weiwei Lin và Victor Chang đã đưa ra mô hình ELSTM bằng làm giảm chiều dữ liệu của chỉ số tỗng hợp và các mã cổ phiếu bằng word vector, rồi đưa ra dự đoán bằng LSTM. Kết quả đạt được chỉ số MSE = 0.017, tốt hơn so với LSTM truyền thống

\subsection{Số feature sử dụng}
Đối với mô hình dự đoán đơn giản, một kích thước đầu vào là đủ.
Tuy nhiên, trong hầu hết các trường hợp, chỉ có một đầu vào không thể đáp ứng
đòi hỏi sự chính xác. Năm 2015, Kai Chen và Yi Zhou [2] có
đã xem xét kỹ ảnh hưởng của số lượng đầu vào đến độ chính xác
trong dự đoán chứng khoán. Từ công việc của họ, người ta có thể thấy rằng
kích thước đầu vào càng nhiều, độ chính xác sẽ đạt được. Với tất cả 5
đầu vào open, close, high, low, value có để đưa ra dự đoán. Ngoài ra, các hiệu ứng của số lượng đầu vào được trình bày bởi Ryo Akita [7], khi sử dụng thông tin bằng số và văn bản cho
dự đoán cổ phiếu. Kết quả là tương tự, càng nhiều độ mờ đầu vào, nhiều yếu tố sẽ được xem xét, cuối cùng sẽ
dẫn đến kết quả tốt hơn. Tuy vậy trong thực tế có một số đầu vào không có đầy đủ tất cả thông tin cho tất cả các phiên, gây rối loạn cho mô hình. Vì vậy chúng tôi chỉ sử dũng 4 cột đầu vào mà có sự đảm bảo về độ đầy đủ của dữ liệu là open, close, high, low.
\subsection{Tối ưu hoá bayesian}
Thuật toán tối ưu hóa Bayes là một thuật toán tối ưu hóa cho một hàm hộp đen nhiễu và có chi phí cao [17] [5]. Thuật toán tối ưu hóa Bayes là chủ yếu
phụ thuộc vào mô hình probilistic xác định phân phối trong
chức năng mục tiêu từ không gian đầu vào để mục tiêu quan tâm.
Đóng góp gần đây cho chức năng tối ưu hóa Bayes
bao gồm các kết quả lý thuyết  [9], đa nhiệm và chuyển giao
tối ưu hóa [4] và ứng dụng vào các tác vụ khác nhau như bộ cảm biến
lựa chọn [10]. Hiệu suất tối ưu hóa bayesian thuật toán trong lựa chọn siêu tham số là ngoài cả mong đợi.	
\section{Mô hình B-LSTM}
Đối với LSTM, có rất ít tham số có thể được thay đổi. Tuy nhiên ta có thể kết hợp với các lớp LSTM lại với nhau và thêm cái các lớp khác để tạo ra nhiều tham số mới. Đối với bài báo này, các siêu tham số tìm kiếm là: số lớp LSTM d, số neuron của các lớp (con số này là bội số của 32), số ngày đầu vào để có thể đưa ra dự đoán cho ngày tiếp theo. Nếu số lớp neuron quá nhiều dẫn đến việc thông tin bị phân tích quá mức, vừa mất thời gian huấn luyện, vừa cho các quả bị overfitting. Nếu quá ít, thì thông tin phân tích không đủ, kết quả đưa ra bị hạn chế về độ chính xác. Số lượng neuron cũng tương tự vậy. Ví LSTM có số ngày đầu vào là cố định nên ta phải tìm số ngày cụ thể đó. Ngoài ra trước khi cho ra kết quả, chúng tôi cho dữ liệu đi qua một lớp dropout để hạn chế tình trạng overfitting, chỉ số dropout có thể lên tới 0.7 để hạn chế việc dữ liệu giao động mạnh ảnh hưởng tới kết quả dự đoán. Việc lạm dụng dropout, luôn chèn lớp dropout vào giữa 2 lớp LSTM có thể dần đến việc mất mát thông tin, mô hình không học được gì, dẫn đến việc mô hình không quan tâm đến đầu vào, mọi kết quả đầu ra đều giống nhau. Vì vậy chúng tối chỉ chèn 1 lớp dropout vào giữa 2 lớp ẩn và lớp output.
\section{Tìm kiếm siêu tham số}
Việc tìm một hàm số có thể biểu diễn đúng tất cả những điểm giá trị trong tập dữ liệu thị trường chứng khoán vừa mất rất nhiều thời gian, tài nguyên vừa có khả năng cao xảy ra trường hợp overfitting. Thay vào đó, yêu cầu đặt ra là một hàm số biểu diễn càng giống với sự phân bố dữ liệu càng tốt, việc này đòi hỏi ít công sức hơn.
Ý tưởng của Bayesian Optimization là dựa trên những bộ tham số trước đó để có thể ra quyết định đến bộ tham số tiếp theo, từ đó hàm số sẽ càng được tối ưu hóa theo thời gian.
Viết theo toán học, việc dựa trên một tập dữ liệu có trước (x) để xây dựng một hàm số với các bộ tham số (y) mô tả bộ dữ liệu đó. Đây là một hàm có điều kiện.
P(y|x)
Như đã nói, việc tối ưu xác suất trong thực tế tốn rất nhiều công sức tính toán, thay vào đó, dựa trên quy tắc Bayes, P(x|y) có thể được tính toán bởi:
P(y|x) = P(x|y)*P(y)/P(x)
Tham khảo quy tắc Bayes tại link:
https://machinelearningcoban.com/2017/07/09/prob/#-quy-tac-bayes

Với x không phục thuộc vào tham số (nên giá trị p(x) không thay đổi khi bộ tham số y thay đổi), Gía trị P(y|x) đạt giá trị lớn nhất khi tích P(x|y)*P(y) đạt gía trị lớn nhất. Việc tối ưu vế phải sẽ dễ dàng hơn, tốt ít tài nguyên hơn vế trái của biểu thức.

Lưu ý: Ta có các cách gọi:
P(x|y): likelihood
P(y|x): posteriori, hoặc surrogate
P(y) : prior

Prior có thể hiểu là gía trị “cảm quan” của dữ liệu. Ví dụ, khi tung một đồng xu đồng chất, xác suất để nhận được mặt sấp-ngửa đều là 1/2. Khi gieo một con xúc sắc, giá trị xác suất nhận được mỗi mặt là 1/6. Đây là giá trị xác suất “cảm quan” nhận được ngay cả khi chưa tung đồng xu hay gieo xúc sắc lần nào. Giá trị này được nhận từ những thông tin của chính dữ liệu mang lại
Có thể hiểu rằng, thay vì dựa trên dữ liệu để tìm ra ngay một hàm số mô tả chính xác bộ dữ liệu đó, ta giả định các bộ dữ liệu một giả thiết biết trước (hàm phân bố dữ liệu, các giá trị được tính toán trên lý thuyết) (prior) và tiến hành tính toán hàm điều kiện likelihood và posterior. Hàm posterior sẽ càng gần với hàm cần tìm sau mỗi vòng lặp cập nhật các tham số.

Nội dung thuật toán:
1. Choose some prior measure over the space of possible objectives f: Chọn một số điểm theo cảm quan (có thể bằng cách tính toán, kinh nghiệm) là tiệm cận với các điểm nằm trên hàm phân bố dữ liệu cần tìm.
2. Combine prior and the likelihood to get a posterior measure over the objective given some observations. : Kết hợp giá trị prior và likelihood để tính toán posterior. Posterior thường sử dụng phương pháp Gaussian Process
3. Use the posterior to decide where to take the next evaluation according to some acquisition/loss function. Thông qua việc tối ưu một hàm acquisition để quyết định vị trí điểm dữ liệu sẽ được đánh giá tiếp theo và cập nhật các tham số cho hàm posterior
4. Augment the data.: thêm điểm dữ liệu 
Iterate between 2 and 4 until the evaluation budget is over : vòng lặp từ bước 2 đến bước 4 kết thúc khi hết tài nguyên (hết số vòng lặp, hết thời gian quy định, ...)

Minh hoạ thuật toán bằng hình: 
http://krasserm.github.io/2018/03/21/bayesian-optimization/

Trong thư viện Python scikit-learn, có thể sử dụng GaussianProcess bằng các gọi hàm GaussianProcessRegressor




\section{Kết quả thực nghiệm}
\subsection{Đầu vào mô hình}
Chúng tôi lấy thông tin của Shanghai A-share Composite Index từ ngày 2006/1/1- 2016/10/19 trên trang web money.163.com. Chúng tối lấy dữ liệu 2 năm cuối để tìm các siêu tham số. Sau đó chúng chia tôi chia dữ liệu thành 70 \% trong số đó là tập huấn luyện, 10\% là tập xác thực và 20\% là tập thử nghiệm. dữ liệu được chuẩn hóa trên [0, 1].
Mô tả cụ thể như sau:
lớn nhất nhỏ nhất

Tối đa là giá trị tối đa của 80\% tập huấn luyện và xác thực và min là mức tối thiểu của ủa 80\% tập huấn luyện và xác thực, sao cho x là
ánh xạ tới [0, 1].
\subsection{Đầu ra mô hình}
Để phân tích hiệu quả của mô hình nhiều hơn
một cách toàn diện, bài viết này đánh giá
hiệu suất của mô hình cho giá cổ phiếu ngắn hạn
dự báo từ hai khía cạnh qua chỉ số cụ thể là MSE trên nguồn dữ liệu đã được ánh xạ từ [0,1].
\subsection{Kết quả}
Để phân tích hiệu quả của mô hình nhiều hơn
một cách toàn diện, bài viết này đánh giá
hiệu suất của mô hình cho giá cổ phiếu ngắn hạn
dự báo từ hai khía cạnh qua chỉ số cụ thể là MSE trên nguồn dữ liệu đã được ánh xạ từ [0,1].
Các mô hình ELSTM, DBN, DBN-MLP và
Các lớp tri giác đa lớp (MLP) tương ứng
Dự đoán thị trường chứng khoán dựa trên mạng lưới thần kinh bộ nhớ ngắn hạn sâu dài thử nghiệm trên hai bộ dữ liệu A share composite index . Kết quả được hiển thị trong hình 5-6.

Có thể thấy, bằng việc chỉ cần tối ưu hoá được bộ tham số, kết quả đã được tối ưu nhiều lần. BLSTM cho chỉ số MSE thấp hơn những 10 lần, hình dạng của biểu đồ dự đoán cũng gần tương tự thực tế, có thể làm cơ sở cho việc đưa ra phán đoán mua bán hoặc giữ giá cổ phiếu

So sánh với E-LSTM
Before you begin to format your paper, first write and save the content as a 
separate text file. Complete all content and organizational editing before 
formatting. Please note sections \ref{AA}--\ref{SCM} below for more information on 
proofreading, spelling and grammar.

Keep your text and graphic files separate until after the text has been 
formatted and styled. Do not number text heads---{\LaTeX} will do that 
for you.

\subsection{Abbreviations and Acronyms}\label{AA}
Define abbreviations and acronyms the first time they are used in the text, 
even after they have been defined in the abstract. Abbreviations such as 
IEEE, SI, MKS, CGS, ac, dc, and rms do not have to be defined. Do not use 
abbreviations in the title or heads unless they are unavoidable.

\subsection{Units}
\begin{itemize}
\item Use either SI (MKS) or CGS as primary units. (SI units are encouraged.) English units may be used as secondary units (in parentheses). An exception would be the use of English units as identifiers in trade, such as ``3.5-inch disk drive''.
\item Avoid combining SI and CGS units, such as current in amperes and magnetic field in oersteds. This often leads to confusion because equations do not balance dimensionally. If you must use mixed units, clearly state the units for each quantity that you use in an equation.
\item Do not mix complete spellings and abbreviations of units: ``Wb/m\textsuperscript{2}'' or ``webers per square meter'', not ``webers/m\textsuperscript{2}''. Spell out units when they appear in text: ``. . . a few henries'', not ``. . . a few H''.
\item Use a zero before decimal points: ``0.25'', not ``.25''. Use ``cm\textsuperscript{3}'', not ``cc''.)
\end{itemize}

\subsection{Equations}
Number equations consecutively. To make your 
equations more compact, you may use the solidus (~/~), the exp function, or 
appropriate exponents. Italicize Roman symbols for quantities and variables, 
but not Greek symbols. Use a long dash rather than a hyphen for a minus 
sign. Punctuate equations with commas or periods when they are part of a 
sentence, as in:
\begin{equation}
a+b=\gamma\label{eq}
\end{equation}

Be sure that the 
symbols in your equation have been defined before or immediately following 
the equation. Use ``\eqref{eq}'', not ``Eq.~\eqref{eq}'' or ``equation \eqref{eq}'', except at 
the beginning of a sentence: ``Equation \eqref{eq} is . . .''

\subsection{\LaTeX-Specific Advice}

Please use ``soft'' (e.g., \verb|\eqref{Eq}|) cross references instead
of ``hard'' references (e.g., \verb|(1)|). That will make it possible
to combine sections, add equations, or change the order of figures or
citations without having to go through the file line by line.

Please don't use the \verb|{eqnarray}| equation environment. Use
\verb|{align}| or \verb|{IEEEeqnarray}| instead. The \verb|{eqnarray}|
environment leaves unsightly spaces around relation symbols.

Please note that the \verb|{subequations}| environment in {\LaTeX}
will increment the main equation counter even when there are no
equation numbers displayed. If you forget that, you might write an
article in which the equation numbers skip from (17) to (20), causing
the copy editors to wonder if you've discovered a new method of
counting.

{\BibTeX} does not work by magic. It doesn't get the bibliographic
data from thin air but from .bib files. If you use {\BibTeX} to produce a
bibliography you must send the .bib files. 

{\LaTeX} can't read your mind. If you assign the same label to a
subsubsection and a table, you might find that Table I has been cross
referenced as Table IV-B3. 

{\LaTeX} does not have precognitive abilities. If you put a
\verb|\label| command before the command that updates the counter it's
supposed to be using, the label will pick up the last counter to be
cross referenced instead. In particular, a \verb|\label| command
should not go before the caption of a figure or a table.

Do not use \verb|\nonumber| inside the \verb|{array}| environment. It
will not stop equation numbers inside \verb|{array}| (there won't be
any anyway) and it might stop a wanted equation number in the
surrounding equation.

\subsection{Some Common Mistakes}\label{SCM}
\begin{itemize}
\item The word ``data'' is plural, not singular.
\item The subscript for the permeability of vacuum $\mu_{0}$, and other common scientific constants, is zero with subscript formatting, not a lowercase letter ``o''.
\item In American English, commas, semicolons, periods, question and exclamation marks are located within quotation marks only when a complete thought or name is cited, such as a title or full quotation. When quotation marks are used, instead of a bold or italic typeface, to highlight a word or phrase, punctuation should appear outside of the quotation marks. A parenthetical phrase or statement at the end of a sentence is punctuated outside of the closing parenthesis (like this). (A parenthetical sentence is punctuated within the parentheses.)
\item A graph within a graph is an ``inset'', not an ``insert''. The word alternatively is preferred to the word ``alternately'' (unless you really mean something that alternates).
\item Do not use the word ``essentially'' to mean ``approximately'' or ``effectively''.
\item In your paper title, if the words ``that uses'' can accurately replace the word ``using'', capitalize the ``u''; if not, keep using lower-cased.
\item Be aware of the different meanings of the homophones ``affect'' and ``effect'', ``complement'' and ``compliment'', ``discreet'' and ``discrete'', ``principal'' and ``principle''.
\item Do not confuse ``imply'' and ``infer''.
\item The prefix ``non'' is not a word; it should be joined to the word it modifies, usually without a hyphen.
\item There is no period after the ``et'' in the Latin abbreviation ``et al.''.
\item The abbreviation ``i.e.'' means ``that is'', and the abbreviation ``e.g.'' means ``for example''.
\end{itemize}
An excellent style manual for science writers is \cite{b7}.

\subsection{Figures and Tables}
\paragraph{Positioning Figures and Tables} Place figures and tables at the top and 
bottom of columns. Avoid placing them in the middle of columns. Large 
figures and tables may span across both columns. Figure captions should be 
below the figures; table heads should appear above the tables. Insert 
figures and tables after they are cited in the text. Use the abbreviation 
``Fig.~\ref{fig}'', even at the beginning of a sentence.

\begin{table}[htbp]
\caption{Table Type Styles}
\begin{center}
\begin{tabular}{|c|c|c|c|}
\hline
\textbf{Table}&\multicolumn{3}{|c|}{\textbf{Table Column Head}} \\
\cline{2-4} 
\textbf{Head} & \textbf{\textit{Table column subhead}}& \textbf{\textit{Subhead}}& \textbf{\textit{Subhead}} \\
\hline
copy& More table copy$^{\mathrm{a}}$& &  \\
\hline
\multicolumn{4}{l}{$^{\mathrm{a}}$Sample of a Table footnote.}
\end{tabular}
\label{tab1}
\end{center}
\end{table}

\begin{figure}[htbp]
\centerline{\includegraphics{fig1.png}}
\caption{Example of a figure caption.}
\label{fig}
\end{figure}

Figure Labels: Use 8 point Times New Roman for Figure labels. Use words 
rather than symbols or abbreviations when writing Figure axis labels to 
avoid confusing the reader. As an example, write the quantity 
``Magnetization'', or ``Magnetization, M'', not just ``M''. If including 
units in the label, present them within parentheses. Do not label axes only 
with units. In the example, write ``Magnetization (A/m)'' or ``Magnetization 
\{A[m(1)]\}'', not just ``A/m''. Do not label axes with a ratio of 
quantities and units. For example, write ``Temperature (K)'', not 
``Temperature/K''.

\section*{References}

Please number citations consecutively within brackets \cite{b1}. The 
sentence punctuation follows the bracket \cite{b2}. Refer simply to the reference 
number, as in \cite{b3}---do not use ``Ref. \cite{b3}'' or ``reference \cite{b3}'' except at 
the beginning of a sentence: ``Reference \cite{b3} was the first $\ldots$''

Number footnotes separately in superscripts. Place the actual footnote at 
the bottom of the column in which it was cited. Do not put footnotes in the 
abstract or reference list. Use letters for table footnotes.

Unless there are six authors or more give all authors' names; do not use 
``et al.''. Papers that have not been published, even if they have been 
submitted for publication, should be cited as ``unpublished'' \cite{b4}. Papers 
that have been accepted for publication should be cited as ``in press'' \cite{b5}. 
Capitalize only the first word in a paper title, except for proper nouns and 
element symbols.

For papers published in translation journals, please give the English 
citation first, followed by the original foreign-language citation \cite{b6}.

\begin{thebibliography}{00}
\bibitem{b1} G. Eason, B. Noble, and I. N. Sneddon, ``On certain integrals of Lipschitz-Hankel type involving products of Bessel functions,'' Phil. Trans. Roy. Soc. London, vol. A247, pp. 529--551, April 1955.
\bibitem{b2} J. Clerk Maxwell, A Treatise on Electricity and Magnetism, 3rd ed., vol. 2. Oxford: Clarendon, 1892, pp.68--73.
\bibitem{b3} I. S. Jacobs and C. P. Bean, ``Fine particles, thin films and exchange anisotropy,'' in Magnetism, vol. III, G. T. Rado and H. Suhl, Eds. New York: Academic, 1963, pp. 271--350.
\bibitem{b4} K. Elissa, ``Title of paper if known,'' unpublished.
\bibitem{b5} R. Nicole, ``Title of paper with only first word capitalized,'' J. Name Stand. Abbrev., in press.
\bibitem{b6} Y. Yorozu, M. Hirano, K. Oka, and Y. Tagawa, ``Electron spectroscopy studies on magneto-optical media and plastic substrate interface,'' IEEE Transl. J. Magn. Japan, vol. 2, pp. 740--741, August 1987 [Digests 9th Annual Conf. Magnetics Japan, p. 301, 1982].
\bibitem{b7} M. Young, The Technical Writer's Handbook. Mill Valley, CA: University Science, 1989.
\end{thebibliography}
\vspace{12pt}
\color{red}
IEEE conference templates contain guidance text for composing and formatting conference papers. Please ensure that all template text is removed from your conference paper prior to submission to the conference. Failure to remove the template text from your paper may result in your paper not being published.

\end{document}
